% Test of where a margin note that doesn't fit on the page of its callout goes, and
% what that does to a heading and to a paragraph's first and last lines.  A note
% fills the margin of the first page, so the note NOTEY doesn't fit there, and a
% rule leaves \testroom lines of the page for the text of the case.  See
% test-noteplace.sh, which runs each case and checks the outcome from the PDF, the
% labels and the warnings only (test-noteplace.py), whatever the mechanism.  The
% cases are those of test-yieldkeep (fix-68), which checked marginfix's log lines.
% NOTEY and its callout carry the labels note-y and callout-y.
\ifdefined\testclass\else\def\testclass{memoir}\fi
\ifdefined\testclassopt\else\def\testclassopt{a4paper}\fi
\documentclass[\testclassopt]{\testclass}
\ifdefined\testoptions\else\def\testoptions{}\fi
\ifdefined\testcase\else\def\testcase{}\fi
\ifdefined\testtall\else\def\testtall{.97}\fi
\expandafter\PassOptionsToPackage\expandafter{\testoptions}{didactic}
\usepackage{didactic}
\usepackage[noamsthm,notheorems]{beamerarticle}
\ifdefined\testpre\testpre\fi
\newcommand\yieldnote[1][]{\label{callout-y}\marginpar{\label{note-y}NOTEY#1}}
% Where the text column is, for test-noteplace.py (in bp, as pdftotext gives):
% from the left edge of odd and of even pages, and its width.
\makeatletter
\AtBeginDocument{\typeout{noteplace geometry:
  odd=\strip@pt\dimexpr(1in+\hoffset+\oddsidemargin)*7200/7227\relax\space
  even=\strip@pt\dimexpr(1in+\hoffset+\evensidemargin)*7200/7227\relax\space
  width=\strip@pt\dimexpr\textwidth*7200/7227\relax\space
  top=\strip@pt\dimexpr(1in+\voffset+\topmargin+\headheight+\headsep)*7200/7227\relax\space
  height=\strip@pt\dimexpr\textheight*7200/7227\relax}}
\makeatother
\newcommand\filler{Some words to fill up the rest of this line of the text and
  also quite a few of the next one, as the words of a paragraph do when it goes
  on for a while, and there is more to say about the matter at hand here.}
% Fills up the line after FIRSTLINE, so that what follows is on the second line.
\newcommand\oneline{Some words to fill up the rest of this line of the text and
  also quite a few}
\def\heading{heading}\def\stacked{stacked}\def\second{second}\def\vertical{vertical}
\def\runin{runin}\def\runintwo{runintwo}\def\club{club}\def\widow{widow}
\def\middle{middle}
\def\theoremcase{theorem}\def\headlist{headlist}\def\withother{withother}
\def\tallheading{tallheading}\def\tallclub{tallclub}\def\tallwidow{tallwidow}
\def\tallmiddle{tallmiddle}\def\split{split}\def\headthm{headthm}
% As club and widow, with a document that forbids the lone line.
\def\clubforbid{clubforbid}\def\widowforbid{widowforbid}
\ifx\testcase\clubforbid
  \let\testcase\club
  \AtBeginDocument{\clubpenalty=10000 \csname @clubpenalty\endcsname=10000 }
\fi
\ifx\testcase\widowforbid
  \let\testcase\widow
  \AtBeginDocument{\widowpenalty=10000 }
\fi
\newcommand\tall{\rule{0pt}{4\baselineskip}}
\newtheorem{thmx}{THMX}
\begin{document}
\ifx\testcase\split
  % The rest of a split note has a box with no key of a callout: it must not
  % stop us.  A long instructor's note is in the way of the note after it.
  \noindent x\ltnote{SPLITNOTE \filler\ \filler\ \filler\ \filler\ \filler\ \filler
    \filler\ \filler\ \filler\ \filler\ \filler\ \filler\ \filler\ \filler\ \filler
    \filler\ \filler\ \filler\ \filler\ \filler\ \filler\ \filler\ \filler\ \filler
    \filler\ \filler\ \filler\ \filler\ \filler\ \filler\ \filler\ \filler\ \filler}
  text
  \par\noindent\rule{1cm}{.7\textheight}\par
  \noindent x\marginpar{ORDINARY} text
  \par\noindent\rule{1cm}{.1\textheight}\par
\else\ifdefined\testfirst
  \testfirst
\else
\noindent x\marginpar{TALL\par\rule{1pt}{\testtall\textheight}} text
\ifdefined\testfootnotefirst\footnote{FOOTA}\fi
\par\noindent\rule{1cm}{\dimexpr\textheight-\testroom\baselineskip-1\baselineskip}\par
\fi\fi
\ifx\testcase\heading
  \section{HEADINGX}
  FIRSTLINE\yieldnote{} \filler\ \filler
\fi
\ifx\testcase\stacked
  \section{HEADINGX}
  \subsection{SUBX}
  FIRSTLINE\yieldnote{} \filler\ \filler
\fi
\ifx\testcase\second
  \section{HEADINGX}
  FIRSTLINE.

  SECONDPARA\yieldnote{} \filler\ \filler
\fi
\ifx\testcase\vertical
  \section{HEADINGX}
  \yieldnote{}
  FIRSTLINE \filler\ \filler
\fi
\ifdefined\testbody\testbody\fi
\ifx\testcase\withother
  \section{HEADINGX}
  FIRSTLINE\yieldnote{} \oneline\ more words\marginpar{NOTEZ} \filler
\fi
% A tall line must not change the count of lines around the break.
\ifx\testcase\tallheading
  \section{HEADINGX}
  \tall FIRSTLINE\yieldnote{} \filler\ \filler
\fi
\ifx\testcase\tallclub
  FIRSTLINE \oneline\ more words\yieldnote{} \filler\tall\ \filler
\fi
\ifx\testcase\tallwidow
  FIRSTLINE \filler\ Now the last line of the paragraph, with its
  note\yieldnote{}\rule{0pt}{3\baselineskip}.
\fi
\ifx\testcase\tallmiddle
  FIRSTLINE \filler\ \filler\yieldnote{} \filler\tall\ \filler\tall\ \filler
\fi
\ifx\testcase\headthm
  \section{HEADINGX}
  \begin{thmx}
  FIRSTLINE\yieldnote{} \filler\ \filler
  \end{thmx}
\fi
\ifx\testcase\headlist
  \section{HEADINGX}
  \begin{itemize}
  \item FIRSTLINE\yieldnote{} \filler\ \filler
  \end{itemize}
\fi
\ifx\testcase\theoremcase
  \begin{thmx}
  FIRSTLINE\yieldnote{} \filler\ \filler
  \end{thmx}
\fi
\ifx\testcase\runin
  \paragraph{HEADINGX}
  FIRSTLINE\yieldnote{} \filler\ \filler
\fi
\ifx\testcase\runintwo
  \paragraph{HEADINGX}
  FIRSTLINE \oneline\ more words\yieldnote{} \filler
\fi
\ifx\testcase\club
  FIRSTLINE \oneline\ more words\yieldnote{} \filler\ \filler
\fi
\ifx\testcase\widow
  FIRSTLINE \filler\ Now the last line of the paragraph, with its
  note\yieldnote{}.
\fi
\ifx\testcase\middle
  FIRSTLINE \filler\ \filler\yieldnote{} \filler\ \filler
\fi

\filler\ \filler\ \filler

\filler\ \filler\ \filler

\newpage
LASTPAGE
\end{document}
